% @LITE_DESC: IEEE conference/journal paper template with two-column layout
% @LITE_SCENE: IEEE paper submissions, computer science conferences, journal articles
% @LITE_TAGS: latex, ieee, paper, conference, journal, two-column, academic

\documentclass[conference]{IEEEtran}

\usepackage{cite}
\usepackage{amsmath, amssymb, amsfonts}
\usepackage{graphicx}
\usepackage{textcomp}
\usepackage{booktabs}
\usepackage{hyperref}

\begin{document}

\title{Your IEEE Paper Title\\
{\footnotesize \textsuperscript{*}Note: Sub-titles are not allowed}
}

\author{\IEEEauthorblockN{First Author}
\IEEEauthorblockA{Affiliation\\
City, Country\\
email@example.com}
\and
\IEEEauthorblockN{Second Author}
\IEEEauthorblockA{Affiliation\\
City, Country\\
email2@example.com}
}

\maketitle

\begin{abstract}
Write a concise abstract (150--250 words) covering the problem, approach, key results, and conclusions.
\end{abstract}

\begin{IEEEkeywords}
keyword1, keyword2, keyword3, keyword4
\end{IEEEkeywords}

\section{Introduction}
\label{sec:intro}
State the problem and contributions. Reference prior work~\cite{ref1}.

\section{Related Work}
\label{sec:related}
Review of prior approaches and their limitations.

\section{Proposed Method}
\label{sec:method}

\subsection{Problem Formulation}
Define the optimization problem:
\begin{equation}
    \min_{\mathbf{x}} f(\mathbf{x}) \quad \text{s.t.} \quad g_i(\mathbf{x}) \leq 0
\end{equation}

\subsection{Algorithm}
Describe the algorithm step by step.

\section{Experiments}
\label{sec:experiments}

\subsection{Setup}
Describe datasets, baselines, and evaluation metrics.

\subsection{Results}
Present quantitative results in Table~\ref{tab:results}.

\begin{table}[htbp]
\caption{Performance comparison on benchmark.}
\label{tab:results}
\centering
\begin{tabular}{lccc}
\toprule
Method & Precision & Recall & F1 \\
\midrule
Baseline A & 0.82 & 0.78 & 0.80 \\
Baseline B & 0.85 & 0.83 & 0.84 \\
\textbf{Ours} & \textbf{0.91} & \textbf{0.89} & \textbf{0.90} \\
\bottomrule
\end{tabular}
\end{table}

\section{Conclusion}
\label{sec:conclusion}
Summarize contributions and discuss limitations.

\bibliographystyle{IEEEtran}
\begin{thebibliography}{10}
\bibitem{ref1} A.~Author, ``Paper title,'' \textit{Journal Name}, vol.~1, pp.~1--10, 2024.
\end{thebibliography}

\end{document}
